\paragraph{Supplement to Finding A52: the norm lower bound and the trace boundary.}
In a compatible cyclic $\ell$-power tower as above, write
$\delta_r=[Q_r-P_r]$, let $b_r=e_{\ell^r}\delta_r$ be its actual
rational primitive projection, and put
\[
 \mathcal O=\mathbf Z_p[Z]/\Phi_{\ell^r}(Z),\qquad
 R=\mathcal O^\tau,\qquad f=\operatorname{ord}(\tau),\qquad
 h=\varphi(\ell^r)/f,\qquad W=\#(Rb_r).
\tag{TN.1}
\]
The rings here are the full finite-etale product algebras, and the
module is the actual primitive torsion module. Flatness in the
descent isomorphism of A45 applies to the inclusion
$Rb_r\subset P_{\ell^r}(k)[p^\infty]$ and gives
\[
                       \mathcal O b_r\simeq
                       \mathcal O\otimes_R Rb_r,
                   \qquad \#(\mathcal O b_r)=W^f.
\tag{TN.2}
\]
In particular bounded rational generator number does not bound the
size of the geometric deck-orbit module.

If $A\in\mathbf Z[C_{\ell^r}]$ annihilates $\delta_r$ and has nonzero
faithful evaluation, then
\[
                              \|A\|_1\ge W^{1/h}.
\tag{TN.3}
\]
Indeed, $\alpha=A(\zeta_{\ell^r})$ annihilates $\mathcal O b_r$.
This is a cyclic $\mathcal O$-module, and the inclusion of ideals
$(\alpha)\subset\operatorname{Ann}_{\mathcal O}(b_r)$ makes
$\mathcal O/(\alpha)$ surject onto it. The nonzero integer
cyclotomic norm therefore satisfies
\[
 W^f\mid N_{\mathbf Q(\zeta_{\ell^r})/\mathbf Q}(\alpha),
 \qquad
 \left|N_{\mathbf Q(\zeta_{\ell^r})/\mathbf Q}(\alpha)\right|
                   \le\|A\|_1^{\varphi(\ell^r)}.
\tag{TN.4}
\]
The first assertion is a divisibility assertion between powers of
$p$ and a nonzero integer, obtained from the $p$-adic index; the
second uses the triangle inequality at every complex embedding.
Together they prove (TN.3). This argument includes every component
of the product algebra, not only a selected character packet.

There is a sharper exact statement when deck Frobenius is transitive
on all faithful characters, so $R=\mathbf Z_p$. If $b_r$ has exact
order $m=p^u$, (TN.2) gives
$\mathcal O b_r\simeq\mathcal O/m\mathcal O$.
Restrict now to integral operators that vanish on \emph{every proper
conductor} and annihilate $b_r$. Their faithful component is divisible
by $m$; all other components are zero. The cyclotomic product
decomposition of $\mathbf Z_p[C_{\ell^r}]$ consequently shows that
each integral coefficient is divisible by $m$. Such an operator
has augmentation zero. Every nonzero one has norm at least $2m$,
and equality is attained by
\[
                         m(1-\sigma^{\ell^{r-1}}).
\tag{TN.5}
\]
Thus $2m$ is the exact minimum in this conductor-isolating class.
It is not a lower bound for arbitrary operators annihilating only
$b_r$: the integral trace $\Phi_{\ell^r}(\sigma)$ already kills
that primitive class with norm $\ell$.

For the full point difference, this cheap trace has a different
meaning. Represent an integral operator by its unique polynomial of
degree less than $\ell^r$. If its faithful value is zero, monic
division in $\mathbf Z[X]$ gives
\[
 A=\Phi_{\ell^r}(\sigma)B,\qquad
 \deg B<\ell^{r-1},\qquad
 \Phi_{\ell^r}(\sigma)=\sum_{j=0}^{\ell-1}
                                      \sigma^{j\ell^{r-1}}.
\tag{TN.6}
\]
For the cover map $v:D_r\to D_{r-1}$ this is the literal divisor
and class identity
\[
 A(Q_r-P_r)=v^*B(Q_{r-1}-P_{r-1}),\qquad
 A\delta_r=v^*B\delta_{r-1}.
\tag{TN.7}
\]
If the latter class is zero, then $B\delta_{r-1}=0$: pullback is
injective on $p$-primary classes because the norm of a pulled-back class
is multiplication by $\ell$ on the original class. A principal function for (TN.7) is
therefore a pullback of the older principal function, up to a
nonzero constant, over $\bar k$. This does not claim arithmetic
descent of a function when $B$ is not Frobenius-invariant.

More generally, for $r>s$ put $d_0=\ell^{r-s}$ and
\[
                 T_{r/s}=\sum_{j=0}^{d_0-1}\sigma^{j\ell^s}.
\tag{TN.8}
\]
Then $T_{r/s}\delta_r=v^*\delta_s$. If an integral operator has
zero cyclotomic components at every conductor above $\ell^s$,
successive monic division gives exactly
\[
             A=T_{r/s}B,\qquad \deg B<\ell^s,
                         \qquad\|A\|_1=d_0\|B\|_1.
\tag{TN.9}
\]
The coefficient supports of the $d_0$ shifted copies of $B$ are
disjoint, which proves the norm equality. The same norm--pullback
argument makes the principal relation descend to $D_s$. Every
deck translate of its principal function still factors through
$D_s$; hence the kernel of the action on the generated field contains the
kernel of $C_{\ell^r}\to C_{\ell^s}$. One cannot use $\ell^r$ as the faithful
group order in the controlled-field argument.

For completeness, even the formally optimistic larger-level
inequality would be stronger. For an integer $d=\|B\|_1\ge1$,
\[
 \frac{(d_0d-1)^2}{d_0}-(d-1)^2
                    =(d_0-1)d^2-1+\frac1{d_0}>0.
\tag{TN.10}
\]
Thus $\ell^r>(d_0d-1)^2$ implies
$\ell^s>(d-1)^2$. In particular, if the exact order of $\delta_s$
is $m_s$, the relation $m_sT_{r/s}$ of norm $d_0m_s$ merely
pulls back the degree-$m_s$ function at level $s$.

These assertions separate two exhaustive cases for a proposed
annihilator at the top level: a nonzero faithful value obeys the
actual orbit-module norm bound, while a zero faithful value traces
down to an older relation. Neither a bounded rational generator
count nor inflation of a trace relation supplies the missing
uniform estimate on the annihilator indices in (IA.14).
